%% ============================================================================
%% Section: The Causal Query Language
%% Operator -> procedure mapping follows dsl/.../exec/CypherTranslator.java
%% and plugin/.../api/ProcedureNames.java (see CODE-PLAN.md, WP2 and WP4).
%% ============================================================================
\section{The Causal Query Language}\label{sec:lang}

Requirement R1 asks for a declarative interface in which each catalog query is
a single statement that compiles to physical operators over the graph store.
In this section we present the language: its surface
(Section~\ref{sec:lang:surface}), the mapping from each operator to the
physical operator that answers it (Section~\ref{sec:lang:compile},
Table~\ref{tab:operators}), the registry through which a database holds
several causal worlds (Section~\ref{sec:lang:registry}), and how a model
enters the store (Section~\ref{sec:lang:populate}). The physical operators
themselves are the subject of Section~\ref{sec:index}.

\subsection{Design and surface syntax}\label{sec:lang:surface}

The language is a small statement-oriented DSL. A program is a sequence of
statements of three kinds: \textit{definition} statements declare variables
and edges and thereby write the causal model into the graph;
\textit{structural} statements ask the catalog queries of
Section~\ref{sec:bg:catalog}; \textit{quantitative} statements express
interventions and counterfactuals over a parameterized model
(Section~\ref{sec:quant}). Two design choices follow from R1. First, an
operator names a causal concept rather than a traversal: an analyst writes
\lstinline[style=causaldsl]|CONFOUNDERS OF T AND Y| and not a path pattern, so
the statement reads as the question and its compilation is the system's
concern. Second, every statement refers to variables by their canonical
\textsf{id}, so a query is portable across any database holding a model with
those variables.

Definition statements come in a standalone and a block form. A bare
\lstinline[style=causaldsl]|DEFINE CAUSAL VARIABLE| introduces one variable
with its observation status, and a
\lstinline[style=causaldsl]|DEFINE CAUSAL DAG| block declares a model as a
set of edges, optionally carrying a structural coefficient per edge:

\begin{lstlisting}[style=causaldsl]
DEFINE CAUSAL DAG CovidModel {
  covid19_burden    -> interventions
  covid19_burden    -> exposure_sarscov2
  interventions     -> mobility
  mobility          -> exposure_sarscov2
  exposure_sarscov2 -> onset_covid19
  onset_covid19     -> reported_cases
}
\end{lstlisting}

\noindent
Against the model so defined, the structural questions of the running example
are each one statement:

\begin{lstlisting}[style=causaldsl]
CONFOUNDERS OF interventions AND reported_cases
BACKDOOR PATHS BETWEEN interventions AND reported_cases
IDENTIFY EFFECT OF interventions ON reported_cases
  RETURN ADJUSTMENT SETS EXPLAIN
\end{lstlisting}

\noindent
The first returns \textit{covid19\_burden}; the second returns the variables
on backdoor paths; the third decides identifiability by adjustment and, with
\lstinline[style=causaldsl]|RETURN ADJUSTMENT SETS|, lists the minimal
adjustment sets, the optional \lstinline[style=causaldsl]|EXPLAIN| attaching
the canonical set and the forbidden set as witnesses of the decision. The
reachability and path-role operators
(\lstinline[style=causaldsl]|CAUSES OF|, \lstinline[style=causaldsl]|EFFECTS OF|,
\lstinline[style=causaldsl]|ANCESTORS OF|, \lstinline[style=causaldsl]|DESCENDANTS OF|,
\lstinline[style=causaldsl]|MEDIATORS BETWEEN|, \lstinline[style=causaldsl]|COLLIDERS BETWEEN|)
follow the same one-statement-per-question shape. The full grammar is given in
ANTLR form in our repository; Table~\ref{tab:operators} lists the operators
that matter for what follows.

\subsection{Compilation to physical operators}\label{sec:lang:compile}

A statement is parsed to an AST and compiled, statement by statement, to one or
more Cypher units executed inside the DBMS. The compiler emits two kinds of
unit. A definition statement compiles to Cypher \textit{writes} that maintain
the schema of Section~\ref{sec:bg:dag}: a \lstinline[style=causaldsl]|DEFINE|
becomes a \texttt{MERGE} of the \textsf{Variable} nodes and \textsf{CAUSES}
edges it introduces (one batched \texttt{UNWIND} for the edge set), after
which the transaction listener of Section~\ref{sec:incr} brings the compiled
snapshot up to date. A structural statement compiles to a \textit{read}: a
\texttt{CALL} of the \texttt{causal.*} procedure that answers it, which
consults the snapshot rather than the stored graph. This two-target
compilation is what keeps the model in the database: the language holds no
copy of the graph, and the same snapshot serves every reader.

\begin{table*}[t]
  \centering
  \caption{Language operators, their graphical semantics, and the physical
  operator each compiles to. Reachability and path-role operators map to
  \texttt{causal.*} procedures over the compiled snapshot
  (Section~\ref{sec:index}); d-separation is exposed as \texttt{causal.dsep}
  and is the inner test of identification. Definition statements compile to
  Cypher writes; world statements resolve in the server-side registry.}
  \label{tab:operators}
  \small
  \setlength{\tabcolsep}{4pt}
  \begin{tabular}{@{}p{0.30\linewidth} p{0.34\linewidth} p{0.30\linewidth}@{}}
    \toprule
    \textbf{DSL statement} & \textbf{Graphical meaning} & \textbf{Physical operator} \\
    \midrule
    \multicolumn{3}{@{}l}{\textit{Definition and worlds}}\\
    \lstinline[style=causaldsl]|DEFINE ... DAG| & write or update the model & \texttt{MERGE} + listener patch \\
    \lstinline[style=causaldsl]|INTERVENE ON X| & $do(X)$: remove edges into $X$ & \texttt{causal.intervene} (derived snapshot) \\
    \lstinline[style=causaldsl]|USE MODEL| / \lstinline[style=causaldsl]|RESET| & select the active world & \texttt{causal.use} \\
    \addlinespace[2pt]
    \multicolumn{3}{@{}l}{\textit{Reachability}}\\
    \lstinline[style=causaldsl]|CAUSES OF X|       & $\pa(X)$ & \texttt{causal.parents} \\
    \lstinline[style=causaldsl]|EFFECTS OF X|      & $\ch(X)$ & \texttt{causal.children} \\
    \lstinline[style=causaldsl]|ANCESTORS OF X|    & $\an(X)$ & \texttt{causal.ancestors} \\
    \lstinline[style=causaldsl]|DESCENDANTS OF X|  & $\de(X)$ & \texttt{causal.descendants} \\
    \addlinespace[2pt]
    \multicolumn{3}{@{}l}{\textit{Path roles}}\\
    \lstinline[style=causaldsl]|MEDIATORS ...|   & $\de(T) \cap \an(Y)$ & \texttt{causal.mediators} \\
    \lstinline[style=causaldsl]|COLLIDERS ...|   & head-to-head meeting points & \texttt{causal.colliders} \\
    \lstinline[style=causaldsl]|CONFOUNDERS ...| & $\an(T) \cap \an(Y)$ & \texttt{causal.confounders} \\
    \lstinline[style=causaldsl]|BACKDOOR PATHS ...| & variables on backdoor paths & \texttt{causal.backdoorVariables} \\
    \addlinespace[2pt]
    \multicolumn{3}{@{}l}{\textit{Independence and identification}}\\
    (inner test) & $X \dsep Y \mid Z$ & \texttt{causal.dsep} \\
    \lstinline[style=causaldsl]|IDENTIFY EFFECT ...| & adjustment criterion, canonical set & \texttt{causal.adjustmentSets} (canonical) \\
    \lstinline[style=causaldsl]|... RETURN ADJUSTMENT SETS| & minimal adjustment sets & \texttt{causal.adjustmentSets} (minimal) \\
    \addlinespace[2pt]
    \multicolumn{3}{@{}l}{\textit{Quantitative}}\\
    \lstinline[style=causaldsl]|INTERVENE ... RETURN EFFECT| & interventional means, total effect & \texttt{causal.linear\allowbreak GaussianInference} \\
    \lstinline[style=causaldsl]|COUNTERFACTUAL ...| & abduction, action, prediction & \texttt{causal.counterfactual} \\
    \bottomrule
  \end{tabular}
\end{table*}

Structural operators map to procedures one-to-one. The reachability operators
compile to \texttt{causal.parents}, \texttt{causal.children},
\texttt{causal.ancestors}, and \texttt{causal.descendants}; the path-role
operators to \texttt{causal.mediators}, \texttt{causal.colliders},
\texttt{causal.confounders}, and \texttt{causal.backdoorVariables};
\lstinline[style=causaldsl]|IDENTIFY EFFECT| compiles to
\texttt{causal.adjustmentSets}, whose mode argument selects the canonical set
(decision only) or the enumeration of minimal sets. D-separation is not a
top-level keyword but the inner test that identification issues; it is exposed
as \texttt{causal.dsep} so that the cost of the traversal-intensive part of
the workload is borne by one physical operator rather than by client-side
orchestration. The quantitative statements compile to
\texttt{causal.linearGaussianInference} and \texttt{causal.counterfactual}
(Section~\ref{sec:quant}).

\subsection{Model registry and causal worlds}\label{sec:lang:registry}

A causal database is rarely a single fixed graph: a curator maintains several
models, and an intervention $do(X)$ defines a derived model, the same DAG with
the edges into $X$ removed, that an analyst wants to name, query, and compare
against the original. The language therefore addresses models by name.
\lstinline[style=causaldsl]|DEFINE MODEL| registers a base model,
\lstinline[style=causaldsl]|USE MODEL| makes a named model (base or derived)
the active one, and \lstinline[style=causaldsl]|RESET| returns to the base
model while leaving derived worlds addressable. Every read procedure accepts
the model name as an optional argument, and the registry of named snapshots
lives in the plugin (Section~\ref{sec:compile}), so that switching worlds is
a name lookup on the server and a derived world is materialized once, by the
graph surgery of Section~\ref{sec:incr:worlds}, and shared by every later
statement that names it.

\subsection{Populating the store}\label{sec:lang:populate}

Definition statements are how a model enters the database, whether typed by
a curator or produced upstream. Everything statistical (structure discovery,
fitting of structural equations, refutation) is by design outside the DBMS:
an accompanying Python pipeline, built on causal-learn and DoWhy, discovers or
accepts a DAG, fits a linear-Gaussian model when requested, and emits a
grammar-valid program of definition statements with a provenance record. The
database receives only the finished structural artifact and performs graph
operations over it. This boundary is what makes the following sections a
database problem rather than a learning problem.
